\chapter{Open Items}\label{r03:ch:open}

\emph{The measurements that were still running when this report was written, the questions
they answer, and the known limitations.}

\section{Measurements still running}\label{r03:sec:open-jobs}

When this report was written (board clock 2026-11-12, about 21:30), four session scripts were
queued on the board, in this order. Each waits for the one before it.

\begin{r03pending}[Pending measurements]
\begin{description}
  \item[1. Halo exchange, end to end] (\ghws{misc/x60-lab/queue_e2e_halo.sh}, running).
        \begin{itemize}
          \item First applies the keep-or-restore rule for the row-pointer change
                (\cref{r03:sec:res-cg}); the change was kept.
          \item Installs the slot-by-slot halo exchange (\cref{r03:sec:wl-halo}) and measures
                all variants in f32 and f64 (logs \texttt{e2e-halo-gcc-\{f32,f64\}.log}).
          \item Then profiles the f32 AA application (\file{perf-halo-aa.txt}) and measures four
                variants with transparent huge pages for the whole application, the scalar builds
                included (\texttt{e2e-halo-thp-gcc-\{f32,f64\}.log}).
        \end{itemize}
  \item[2. Padding between the PDF slices, lab] (\ghws{misc/x60-lab/queue_lab6.sh}): extra
        bytes between the 19 slices (0 to 33\,024), with 4~KiB pages and with huge pages, for the
        adopted pull and AA kernels and two scalar baselines (\file{lab6.log}).
  \item[3. Kernels without output and the split kernel, lab] (\ghws{misc/x60-lab/queue_lab7.sh}),
        3 rounds in one session (\file{lab7.log}):
        \begin{itemize}
          \item the adopted kernels with and without the density/velocity output;
          \item the scalar AA kernel with and without output (generated without CSE, see
                \cref{r03:sec:open-limits});
          \item the hand kernels as references (GCC; Clang with the X60 model);
          \item the generated split kernel over BY, GM, GP, with and without prefetch and output
                (GCC; the main configuration also with Clang and the X60 model).
        \end{itemize}
  \item[4. Kernels without output and the split kernel, end to end]
        (\ghws{misc/x60-lab/queue_e2e_final.sh}).
        \begin{itemize}
          \item Installs the kernels without output and the split variant, and regenerates every
                variant.
          \item Measures all variants with output in every step
                (\texttt{e2e-final-gcc-\{f32,f64\}.log}) and with output in the last step only
                (\texttt{e2e-final-last-gcc-\{f32,f64\}.log}).
          \item Then profiles the f32 split application (\file{perf-final-split.txt}).
        \end{itemize}
\end{description}
\end{r03pending}

These runs fill the pending rows of \cref{r03:tab:waterfall32,r03:tab:waterfall64}.

\section{Questions these runs answer}\label{r03:sec:open-questions}

\begin{itemize}
  \item How much of the halo exchange's 21\,\% of an AA time step does the slot-by-slot copy
        recover?
  \item Do huge pages help the AA application, as they helped the original kernel (+7\,\%,
        \cref{r03:sec:bw}), or hurt it, as in the lab (\cref{r03:tab:align})? Does padding
        between the PDF slices remove the penalty?
  \item What does skipping the density/velocity output gain for the adopted kernels, scalar and
        vector?
  \item Does the generated split kernel reach the hand kernel's 164 cycles per cell, and does it
        beat the single-pass AA kernel (192) in the application? Which BY, GM, GP and which
        compiler suit it best?
\end{itemize}

\section{Limitations and follow-ups}\label{r03:sec:open-limits}

\begin{itemize}
  \item \textbf{The split is generated by a script, not by a \ps{} transformation.}
        \code{lbm_split.py} builds ordinary \ps{} kernels and a C++ driver from lbmpy's
        description of the method. It supports SRT with compressible, zero-centred PDFs and AA
        streaming only. A general loop-splitting transformation in \ps{}, or a sweepgen prefab,
        would make it available to other methods and applications.
  \item \textbf{No Clang end-to-end measurement since run A.} Clang's application numbers
        predate AA streaming, the new instruction selection and the scheduling-model flag. Clang
        gains 9 to 31\,\% from the X60 model in the lab, but GCC still builds the faster
        generated AA kernel.
  \item \textbf{The scalar baselines have no software prefetch.} The \ps{} option exists for
        vector loops only. In the lab, a scalar kernel with prefetch but without CSE was 15\,\%
        faster than the scalar kernel with CSE (\cref{r03:tab:pf}); prefetch on the scalar kernel
        with CSE was not measured. A scalar prefetch would make the best scalar baseline
        stricter.
  \item \textbf{The lab's scalar AA kernels have no CSE.} The same naming error
        (\cref{r03:sec:repro-gen}) applies to \file{f32-generic-cse_scalar+aa.cpp} and its
        variants, which the pending lab session measures. Their results compare scalar AA kernels
        \emph{without} CSE; the application's best scalar build has CSE.
  \item \textbf{The GPU path of the halo-exchange rewrite was not run.} The slot-by-slot pack
        infos also serve GPU execution (one kernel launch per slot); no GPU was available to test
        them.
  \item \textbf{The hand kernels use the textbook equilibrium.} Their speed figures are valid,
        but they do not compute lbmpy's model exactly (\cref{r03:sec:wl-split}).
  \item \textbf{Several processes} were measured once (\cref{r03:sec:res-mpi}) and not
        followed up: on one embedded system on a chip, the single-core figures are the result,
        and the many-stream pattern saturates the chip's memory at about 3.6~GB/s anyway.
\end{itemize}
